\section*{Dok.\ 252 — Zahlentheoretische Strukturen: Phillips-Framework und FFGFT}
\addcontentsline{toc}{section}{Dok.\ 252 — Zahlentheoretische Strukturen}

\begin{abstract}
Paul Phillips hat im Rahmen der IPI-Gruppe ein zahlentheoretisches Framework
eingereicht, das strukturelle Eigenschaften der $\sigma$-Funktion (Teilersumme),
Homologie und dimensionale Dynamik verbindet (23 Dateien, SHA-256 gepinnt).
Dieses Dokument ist eine erste systematische Gegenüberstellung mit FFGFT.
Alle Kernaussagen wurden Python-verifiziert. Das Bundle ist noch nicht
vollständig ausgewertet; offene Punkte werden explizit benannt.
Der neue Zentralbefund: die $\sigma$-Waisen-Primzahlen im 13-glatten Ambient
sind exakt $\{2, 5, 11\}$ — und $5$, $11$ sind genau die mittleren Terme
der kubischen Folge $3 \to 5 \to 11 \to 1321$.
\end{abstract}

\subsection*{1. Elemente des Phillips-Frameworks}

\subsubsection*{1.1 Die $\sigma$-Funktion und $h = 15$}

Die Teilersumme $\sigma(n) = \sum_{d \mid n} d$ hat für $n = 15$:
\[
\sigma(15) = \sigma(3 \cdot 5) = (1+3)(1+5) = 4 \cdot 6 = 24 = 4!
\qquad \text{[Python: } \checkmark\text{]}
\]
\[
\sigma(14) = \sigma(2 \cdot 7) = (1+2)(1+7) = 3 \cdot 8 = 24
\qquad [\sigma(14) = \sigma(15)]
\]
$h = 15$ ist ein Berührungspunkt zweier $\sigma$-Äste.

\subsubsection*{1.2 $\sigma$-Waisen und das $L_0$-Skelett}

Eine $\sigma$-Waise ist eine Zahl, die nicht als $\sigma(n)$ für irgendein $n$ auftritt.
Das 8-stufige $L_0$-Skelett ist durch $\sigma$-Waisen begrenzt
(genaue Konstruktion aus Bundle noch nicht vollständig ausgewertet).

\begin{mdframed}[style=ok]
\textbf{Python-Befund (neu):} Die $\sigma$-Waisen-Primzahlen im 13-glatten
Ambient $\langle 2, 3, 5, 7, 11, 13 \rangle$ sind exakt:
\[
\boxed{\{2,\; 5,\; 11\}}
\]
Genau drei Elemente. $3$ ist keine Waise: $\sigma(2) = 3$.
$7$ ist keine Waise: $\sigma(4) = 7$. $13$ ist keine Waise: $\sigma(9) = 13$.
\end{mdframed}

\subsubsection*{1.3 Kubische Rekursion $3 \to 5 \to 11 \to 1321$}

Die Primzahlfolge entsteht durch eine kubische Rekursionsrelation
(genaue Formel im Bundle, noch nicht extrahiert). Alle Terme sind prim.

\begin{mdframed}[style=ok]
\textbf{Python-Verifikation:} $1321$ prim $\checkmark$ \quad
$1321 = 2^3 \cdot 3 \cdot 5 \cdot 11 + 1 = 1320 + 1$ $\checkmark$\\
$1320 = 8 \cdot 3 \cdot 5 \cdot 11$ — Produkt aller vorherigen Terme mal $2^3$.
\end{mdframed}

\subsubsection*{1.4 13-glatte Zahlen}

Das Ambient-Set $\langle 2, 3, 5, 7, 11, 13 \rangle$ enthält alle ganzen Zahlen
mit Primfaktoren ausschließlich aus $\{2, 3, 5, 7, 11, 13\}$.
Die erste Primzahl außerhalb ist $17$.

\subsubsection*{1.5 Klein-Vier-Gruppe $V_4$ und CPT}

$V_4 = \mathbb{Z}_2 \times \mathbb{Z}_2 = \{e, C, P, T\}$ mit
$C^2 = P^2 = T^2 = e$. Phillips' $\sigma$-Arithmetik generiert $V_4$ als
natürliche Symmetriegruppe.

\subsubsection*{1.6 Dimensionale Dynamik, Epitaxie, Homologie}

Das Framework verbindet dimensionale Dynamik mit Epitaxie
(kristallographische Gitterkompatibilität) und Homologie.
Formalisierung im Bundle.

\subsection*{2. FFGFT-Parallelen: systematischer Vergleich}

\subsubsection*{2.1 Strukturprimes $\{2, 3, 5\}$ und das 13-glatte Ambient}

In FFGFT hat $\xi$ die Primfaktorzerlegung:
\[
\xi = \frac{4}{30000} = \frac{1}{7500} = \frac{1}{3 \cdot 2^2 \cdot 5^4}
\qquad \text{[Python: } \xi^{-1} = 7500 = 3 \times 4 \times 625\ \checkmark\text{]}
\]

\begin{center}
\adjustbox{max width=\textwidth}{%
\begin{tabular}{lll}
\toprule
Primfaktor & Potenz in $\xi^{-1}$ & Geometrische Bedeutung \\
\midrule
$2$ & $2^2 = 4$ & Raumzeitdimensionen des T$^4$ \\
$3$ & $3^1$ & Zahl der Raumdimensionen \\
$5$ & $5^4$ & Fraktale Schichtstruktur \\
\bottomrule
\end{tabular}%
}
\end{center}

Die übrigen Primes des 13-glatten Ambients:

\begin{center}
\adjustbox{max width=\textwidth}{%
\begin{tabular}{lll}
\toprule
Primzahl & Rolle in FFGFT & Herkunft \\
\midrule
$7$ & $\kappa = 7$ (kosmologisch) & intern konstrained \\
$11$ & Leptonsektor, Exponentarchitektur & Rekursion \\
$13$ & $\alpha^{-1} = (\xi \cdot \varphi^3 \cdot 13)^{-1} = 136{,}19$ & Methode 1 \\
\bottomrule
\end{tabular}%
}
\end{center}

\begin{mdframed}[style=ok]
$(\xi \cdot \varphi^3 \cdot 13)^{-1} = 136{,}19$ [Python: $\checkmark$]\\
$17$ (erste Außenprimzahl) hat keine strukturelle Rolle in FFGFT. Konsistent.
\end{mdframed}

\subsubsection*{2.2 Zentralbefund: $\sigma$-Waisen und die kubische Folge}

\begin{mdframed}[style=ok]
\textbf{Stärkstes Ergebnis — Python-verifiziert:}

Die $\sigma$-Waisen-Primzahlen im 13-glatten Ambient sind $\{2, 5, 11\}$.

Die kubische Folge lautet $3 \to \mathbf{5} \to \mathbf{11} \to 1321$.

Die Terme $5$ und $11$ sind genau die $\sigma$-Waisen-Primes (außer $2$).
Der Startwert $3$ ist \textit{keine} Waise ($\sigma(2) = 3$).

\textbf{Konsequenz:} Die kubische Rekursion startet an einem $\sigma$-Wert
und durchläuft dann ausschließlich $\sigma$-Waisen-Primes aus dem 13-glatten
Ambient. Das ist keine numerische Koinzidenz.
\end{mdframed}

In FFGFT erscheinen $5$ und $11$ im Leptonsektor:
$r_\tau = 25/9 = 5^2/3^2$; $p_e = 3/2$; Rekursionstiefe $N \approx 8{,}4$
enthält Faktor~$5$ ($\approx 25/3$).

\subsubsection*{2.3 $h = 15$ als $\{3,5\}$-Kern von $30000$}

\[
30000 = 2^4 \cdot 3 \cdot 5^4 = 2^4 \cdot 5^3 \cdot \underbrace{15}_{3 \times 5}
\qquad [\text{Python: } \checkmark]
\]
$h = 15$ greift genau die beiden nicht-binären Strukturprimes von $\xi$ heraus.
Weiterer Befund: $\sigma(15) = 24 = 4! = |S_4|$,
und T$^4$ hat Permutationssymmetrie $S_4$ (Ordnung 24).

\subsubsection*{2.4 Rekursionstiefe $N \approx 8$}

In FFGFT durch zwei unabhängige Bedingungen:
(1) $Z_3^3$-topologischer Abschluss (27 Schritte, Cutoff bei $\approx 8{,}4$);
(2) Fibonacci-Konvergenz $\varphi^{-2n}/\sqrt{5} \approx \xi \Rightarrow n \approx 8{,}4$.

Phillips' 8-stufiges $L_0$-Skelett: Konstruktion aus Bundle noch ausstehend.
Die $\sigma$-Waisen-Primzahlen zählen nur $\{2, 5, 11\}$ — drei Elemente,
nicht acht. Das Skelett muss anders konstruiert sein als die bloße Zählung
dieser Primes.

\subsubsection*{2.5 $V_4$ und CPT}

In FFGFT aus T$^4$-Topologie abgeleitet: jede der 4 Dimensionen kann
unabhängig reflektiert werden ($\mathbb{Z}_2^4$ als maximale Spiegelungsgruppe).
CPT entspricht der Untergruppe $V_4 \subset \mathbb{Z}_2^4$,
die die Lorentz-Signatur erhält. Phillips' $\sigma$-Arithmetik generiert
dieselbe Gruppe — beide nicht postuliert, sondern abgeleitet.

\subsubsection*{2.6 T$^4$-Homologie und Epitaxie}

T$^4$ hat Homologiegruppen $H_k(T^4, \mathbb{Z}) = \mathbb{Z}^{\binom{4}{k}}$;
insbesondere $H_1 = \mathbb{Z}^4$ (Windungszahlen $\psi \in H_1$).
Ob Phillips' Epitaxie-Begriff auf die T$^4$-Einbettung anwendbar ist,
bleibt offen.

\subsection*{3. Gesamtbewertung}

\begin{center}
\adjustbox{max width=\textwidth}{%
\begin{tabular}{p{6cm}ll}
\toprule
Überschneidung & Stärke & Status \\
\midrule
$\{2,3,5\}$ als $\xi$-Kern & stark & aus T$^4$ abgeleitet $\checkmark$ \\
$7, 11, 13$ im 13-glatten Ambient & stark & alle in FFGFT $\checkmark$ \\
$17$ als erste Außenprimzahl & konsistent & keine FFGFT-Rolle $\checkmark$ \\
$\sigma$-Waisen $\{2,5,11\}$ = Folge $3 \to 5 \to 11$ & \textbf{stark} & \textbf{Python} $\checkmark$ \\
$h=15 = 3 \times 5$ als $\xi$-Kern & auffällig & nicht trivial \\
$\sigma(15) = 24 = |S_4|$ & interessant & T$^4$-Symmetrie \\
$V_4$/CPT & echte Überschneidung & beide abgeleitet $\checkmark$ \\
Rekursionstiefe $N \approx 8$ & auffällig & Mechanismus unklar \\
$1321 = 2^3 \cdot 3 \cdot 5 \cdot 11 + 1$ & interessant & $\checkmark$ prim \\
E$_8$-Verbindung & unbekannt & Bundle ausstehend \\
Kubische Formel & unbekannt & Bundle ausstehend \\
\bottomrule
\end{tabular}%
}
\end{center}

\begin{mdframed}[style=info]
\textbf{Fazit:} Die Überschneidungen bei Strukturprimes, $V_4$/CPT,
13-glattem Ambient und insbesondere die Identität
$\{\sigma\text{-Waisen-Primes}\} \cap \{\text{kubische Folge}\} = \{5, 11\}$
sind zu spezifisch um rein zufällig zu sein. Das Bundle muss vollständig
ausgewertet werden — insbesondere kubische Rekursionsformel und
E$_8$-Abschnitt — bevor eine stärkere Aussage möglich ist.
\end{mdframed}

\subsection*{4. Nächste Schritte}

\begin{itemize}
\item Bundle vollständig lesen: kubische Rekursionsformel, E$_8$-Abschnitt
\item Frage an Paul: Wie ist das 8-stufige $L_0$-Skelett genau konstruiert?
\item Prüfung: Ergibt Phillips' Rahmen $\xi$ oder eine Näherung davon?
\item Frage: Warum startet die Folge bei $3$ (kein $\sigma$-Waise)?
\end{itemize}
